The abstract machine of \parencite{jones_2026} organises machine states into
\emph{configurations}: binary trees whose leaves are either basic machines or
failures. A configuration transition selects a non-terminal leaf and steps it,
possibly splitting it into further branches. The theory leaves the choice of
which leaf to step entirely to the implementer:

\begin{quote}
  \emph{``Observe that this is highly nondeterministic: it may be that one or
  more leaves of $\Conf{C}$ may be able to make a basic transition! We leave
  the choice to the implementer, who may use any strategy to explore the tree
  (e.g.~DFS, BFS, parallel).''} \hfill---\parencite{jones_2026}
\end{quote}

\noindent This freedom is justified by the equational theory: the denotational
semantics is based on the lower powerdomain, which is insensitive to the order
in which branches are explored. All search strategies produce the same set of
results for terminating programs. Nevertheless, the choice of strategy has
profound practical consequences: it determines which solutions are found first,
whether all solutions are eventually enumerated, and how much memory is
consumed.

\Gweek{} implements four strategies. We describe each, and then discuss their
trade-offs.

\subsection{Breadth-First Search (\BFS{})}

The default strategy maintains two vectors: the current generation and the next
generation. Each machine in the current generation is stepped to its next branch
point; the resulting branches are placed in the next generation. When the
current generation is exhausted, the vectors are swapped.

\BFS{} explores all machines at depth $d$ before any at depth $d+1$. It is
therefore \emph{complete}: if a solution exists at any finite depth, it will
eventually be found. It is also \emph{fair}: no branch is starved indefinitely.
Its principal limitation is memory: the number of active machines can grow
exponentially in the depth of the search tree.

\subsection{Depth-First Search (\DFS{})}

\DFS{} maintains a stack of machines and always steps the most recently created
branch. It uses memory proportional to the depth of the tree, which is typically
far less than its width.

However, \DFS{} is \emph{incomplete}: if the search tree has an infinite branch,
\DFS{} may follow it forever without discovering solutions on other branches.
Infinite branches arise naturally in FLP when logical variables of type $\NatT$
or $\ListT{A}$ are narrowed---each narrowing of a successor or cons constructor
introduces a fresh variable, creating an infinite chain. \DFS{} is therefore
appropriate only for programs with finite search spaces, or when a single
solution is sufficient and the programmer can guarantee termination.

\subsection{Iterative Deepening (\IDDFS{})}

Iterative deepening combines the completeness of \BFS{} with the memory
efficiency of \DFS{}. The strategy runs \DFS{} with a depth limit $d$ that
starts at 1 and doubles after each complete pass. Solutions found at shallower
depths in previous passes are deduplicated by a hash set of previously seen
outputs.

The overhead of re-exploring shallow nodes is bounded by a constant factor
(since the tree grows exponentially). \IDDFS{} is complete, uses $O(d)$ memory
where $d$ is the depth of the shallowest solution, and is most useful when the
search tree is both deep and wide.

\subsection{Fair Round-Robin (\Fair{})}

The fair strategy maintains a queue of \emph{threads}, where each thread is a
local \DFS{} work stack. A single thread is created initially; new threads are
spawned at branch points. When the machine encounters a branch while executing a
thread, the first alternative continues in the current thread and each remaining
alternative is placed in the queue as a new single-machine thread. Each thread
is allowed a fixed quota (currently $10{,}000$ steps). If a thread exhausts its
quota without completing, its remaining work stack is re-enqueued as a single
entry, and the scheduler moves to the next thread.

The key invariant is that \emph{branch alternatives are distributed across
threads}, so every branch receives its own share of the round-robin schedule.
Within a thread, execution is depth-first, so the local work stack remains
small (proportional to the depth of the sub-tree being explored). To bound
memory, the total number of threads is capped (currently $10{,}000$); when the
cap is reached, new branch alternatives remain in their parent thread and are
explored depth-first. This degrades gracefully: existing threads continue to
receive fair scheduling, while only new branches within a full thread lose
individual time-slicing.

This strategy is both complete and fair: every branch receives a bounded share
of computation, preventing any single branch from starving others. It avoids
the memory blow-up of \BFS{} while retaining completeness, and is particularly
effective for programs that mix finite branching (explicit
\mintinline{haskell}|<>|) with infinite branching (logical variable narrowing),
as the finite branches are fully resolved within their quanta while the
infinite branches are interleaved.

\subsection{Discussion}

The theoretical license to choose any exploration order is powerful, but it
places a practical burden on the programmer. Consider a program that introduces
a logical variable $x : \NatT$ and constrains it by $x \mathbin{\texttt{=:=}}
5$. The constraint is satisfied by a single value, but the narrowing of $x$
produces an infinite family of branches: $x = 0$, $x = 1$, $x = 2$, \ldots
Under \BFS{}, the solution $x = 5$ is found at depth~5. Under \DFS{}, the
machine follows the chain $x = 0, x = \textsf{S}(0), x = \textsf{S}^2(0),
\ldots$ and finds it at step~5 as well---but only because the infinite branch
happens to contain the solution. If the constraint were instead $x
\mathbin{\texttt{=:=}} y$ for another logical variable $y$, \DFS{} would
diverge on the very first narrowing of $x$.

In practice, \BFS{} is the safest default for programs involving logical
variables, as it does not require the programmer to reason about the finiteness
of branches. \DFS{} is fastest for finite search spaces. The fair strategy
offers a practical compromise for mixed workloads, and is the recommended
choice when the structure of the search space is unknown.
